% In total max. 1 Page!
\AIRstudentthesisAbstract{%
%
% Abstract English:
Threading a nut onto a bolt with a multi-fingered hand is a long-horizon, contact-rich assembly task. A single turn is limited by the range of the wrist, so the hand must grasp, turn, release, return, and re-grasp the nut in repeated cycles, and small errors at the boundaries between these phases accumulate over the cycles. This thesis studies how such a cyclic skill can be learned for a Franka Research~3 arm with an Inspire RH56DFX dexterous hand and transferred from simulation to hardware. Learning all phases end to end with one reinforcement-learning policy did not succeed, because a reward that captures both long-horizon progress and precise phase transitions could not be designed. The thesis therefore follows the principle of decomposing during data generation and unifying at deployment. A hybrid teacher combines a privileged reinforcement-learning policy for approach, grasp, rotation, and goal reaching with scripted release-and-return controllers to generate complete multi-cycle demonstrations in a simplified thread simulation. Replaying these demonstrations on the physical robot with sparse human corrections yields a mixed dataset that reduces the sim-to-real gap. The demonstrations are distilled into a single visuomotor student by conditional flow matching: a trajectory Transformer predicts eight-step action chunks from point-cloud features of an aligned RGB-D observation, proprioception, a global trajectory-progress scalar, and a process-phase encoding. The thesis derives the flow-matching objective, describes the student as implemented, and compares offline flow matching with online DAgger variants, deterministic LSTM and Transformer regression heads, point-cloud and image encoders, privileged and visual observations, and two action contracts. In ten-cycle trials the distilled policy completes all cycles on M24, M30, and M36 hexagonal nuts in simulation and on hardware, with mean cyclic tracking errors below $6\degree$, and transfers to unseen round and octagonal profiles, while rectangular profiles expose a limit of the learned re-grasp strategy.
}{%
%
% Zusammenfassung Deutsch:
Das Aufschrauben einer Mutter auf einen Bolzen mit einer mehrfingrigen Hand ist eine langfristige, kontaktreiche Montageaufgabe. Eine einzelne Drehung ist durch den Bewegungsbereich des Handgelenks begrenzt, sodass die Hand die Mutter in wiederholten Zyklen greifen, drehen, loslassen, zurückkehren und erneut greifen muss; kleine Fehler an den Phasengrenzen akkumulieren über die Zyklen. Diese Arbeit untersucht, wie eine solche zyklische Fertigkeit für einen Franka-Research-3-Arm mit einer geschickten Inspire-RH56DFX-Hand gelernt und von der Simulation auf die Hardware übertragen werden kann. Das Ende-zu-Ende-Lernen aller Phasen mit einer einzigen Reinforcement-Learning-Policy gelang nicht, da keine Belohnungsfunktion gefunden wurde, die sowohl langfristigen Fortschritt als auch präzise Phasenübergänge erfasst. Die Arbeit folgt daher dem Prinzip, bei der Datenerzeugung zu zerlegen und beim Einsatz zu vereinen. Ein hybrider Lehrer kombiniert eine privilegierte Reinforcement-Learning-Policy für Annäherung, Griff, Drehung und Zielerreichung mit skriptgesteuerten Loslass- und Rückkehrreglern und erzeugt so vollständige mehrzyklische Demonstrationen in einer vereinfachten Gewindesimulation. Die Wiedergabe dieser Demonstrationen auf dem realen Roboter mit sparsamen menschlichen Korrekturen liefert einen gemischten Datensatz, der die Lücke zwischen Simulation und Realität verringert. Die Demonstrationen werden mittels bedingtem Flow Matching in eine einzige visuomotorische Schüler-Policy destilliert: Ein Trajektorien-Transformer sagt Aktionsblöcke von acht Schritten aus Punktwolkenmerkmalen einer ausgerichteten RGB-D-Beobachtung, Propriozeption, einem globalen Fortschrittsskalar und einer Prozessphasenkodierung vorher. Die Arbeit leitet das Flow-Matching-Ziel her, beschreibt den Schüler wie implementiert und vergleicht offline Flow Matching mit online DAgger-Varianten, deterministischen LSTM- und Transformer-Regressionsköpfen, Punktwolken- und Bildkodierern, privilegierten und visuellen Beobachtungen sowie zwei Aktionsschnittstellen. In Versuchen mit zehn Zyklen vollendet die destillierte Policy alle Zyklen auf M24-, M30- und M36-Sechskantmuttern in Simulation und auf der Hardware mit mittleren zyklischen Nachführfehlern unter $6\degree$ und überträgt sich auf ungesehene runde und achteckige Profile, während rechteckige Profile eine Grenze der gelernten Nachgreifstrategie aufzeigen.
}%
%
%
